\documentclass[ignorenonframetext,]{beamer}
\setbeamertemplate{caption}[numbered]
\setbeamertemplate{caption label separator}{: }
\setbeamercolor{caption name}{fg=normal text.fg}
\beamertemplatenavigationsymbolsempty
\usepackage{lmodern}
\usepackage{amssymb,amsmath}
\usepackage{ifxetex,ifluatex}
\usepackage{fixltx2e} % provides \textsubscript
\ifnum 0\ifxetex 1\fi\ifluatex 1\fi=0 % if pdftex
  \usepackage[T1]{fontenc}
  \usepackage[utf8]{inputenc}
\else % if luatex or xelatex
  \ifxetex
    \usepackage{mathspec}
  \else
    \usepackage{fontspec}
  \fi
  \defaultfontfeatures{Ligatures=TeX,Scale=MatchLowercase}
\fi
% use upquote if available, for straight quotes in verbatim environments
\IfFileExists{upquote.sty}{\usepackage{upquote}}{}
% use microtype if available
\IfFileExists{microtype.sty}{%
\usepackage{microtype}
\UseMicrotypeSet[protrusion]{basicmath} % disable protrusion for tt fonts
}{}
\newif\ifbibliography
\hypersetup{
            pdftitle={Class 9: Bayesian Hierarchical Models},
            pdfauthor={Andrew Parnell},
            pdfborder={0 0 0},
            breaklinks=true}
\urlstyle{same}  % don't use monospace font for urls
\usepackage{color}
\usepackage{fancyvrb}
\newcommand{\VerbBar}{|}
\newcommand{\VERB}{\Verb[commandchars=\\\{\}]}
\DefineVerbatimEnvironment{Highlighting}{Verbatim}{commandchars=\\\{\}}
% Add ',fontsize=\small' for more characters per line
\usepackage{framed}
\definecolor{shadecolor}{RGB}{248,248,248}
\newenvironment{Shaded}{\begin{snugshade}}{\end{snugshade}}
\newcommand{\KeywordTok}[1]{\textcolor[rgb]{0.13,0.29,0.53}{\textbf{#1}}}
\newcommand{\DataTypeTok}[1]{\textcolor[rgb]{0.13,0.29,0.53}{#1}}
\newcommand{\DecValTok}[1]{\textcolor[rgb]{0.00,0.00,0.81}{#1}}
\newcommand{\BaseNTok}[1]{\textcolor[rgb]{0.00,0.00,0.81}{#1}}
\newcommand{\FloatTok}[1]{\textcolor[rgb]{0.00,0.00,0.81}{#1}}
\newcommand{\ConstantTok}[1]{\textcolor[rgb]{0.00,0.00,0.00}{#1}}
\newcommand{\CharTok}[1]{\textcolor[rgb]{0.31,0.60,0.02}{#1}}
\newcommand{\SpecialCharTok}[1]{\textcolor[rgb]{0.00,0.00,0.00}{#1}}
\newcommand{\StringTok}[1]{\textcolor[rgb]{0.31,0.60,0.02}{#1}}
\newcommand{\VerbatimStringTok}[1]{\textcolor[rgb]{0.31,0.60,0.02}{#1}}
\newcommand{\SpecialStringTok}[1]{\textcolor[rgb]{0.31,0.60,0.02}{#1}}
\newcommand{\ImportTok}[1]{#1}
\newcommand{\CommentTok}[1]{\textcolor[rgb]{0.56,0.35,0.01}{\textit{#1}}}
\newcommand{\DocumentationTok}[1]{\textcolor[rgb]{0.56,0.35,0.01}{\textbf{\textit{#1}}}}
\newcommand{\AnnotationTok}[1]{\textcolor[rgb]{0.56,0.35,0.01}{\textbf{\textit{#1}}}}
\newcommand{\CommentVarTok}[1]{\textcolor[rgb]{0.56,0.35,0.01}{\textbf{\textit{#1}}}}
\newcommand{\OtherTok}[1]{\textcolor[rgb]{0.56,0.35,0.01}{#1}}
\newcommand{\FunctionTok}[1]{\textcolor[rgb]{0.00,0.00,0.00}{#1}}
\newcommand{\VariableTok}[1]{\textcolor[rgb]{0.00,0.00,0.00}{#1}}
\newcommand{\ControlFlowTok}[1]{\textcolor[rgb]{0.13,0.29,0.53}{\textbf{#1}}}
\newcommand{\OperatorTok}[1]{\textcolor[rgb]{0.81,0.36,0.00}{\textbf{#1}}}
\newcommand{\BuiltInTok}[1]{#1}
\newcommand{\ExtensionTok}[1]{#1}
\newcommand{\PreprocessorTok}[1]{\textcolor[rgb]{0.56,0.35,0.01}{\textit{#1}}}
\newcommand{\AttributeTok}[1]{\textcolor[rgb]{0.77,0.63,0.00}{#1}}
\newcommand{\RegionMarkerTok}[1]{#1}
\newcommand{\InformationTok}[1]{\textcolor[rgb]{0.56,0.35,0.01}{\textbf{\textit{#1}}}}
\newcommand{\WarningTok}[1]{\textcolor[rgb]{0.56,0.35,0.01}{\textbf{\textit{#1}}}}
\newcommand{\AlertTok}[1]{\textcolor[rgb]{0.94,0.16,0.16}{#1}}
\newcommand{\ErrorTok}[1]{\textcolor[rgb]{0.64,0.00,0.00}{\textbf{#1}}}
\newcommand{\NormalTok}[1]{#1}
\usepackage{graphicx,grffile}
\makeatletter
\def\maxwidth{\ifdim\Gin@nat@width>\linewidth\linewidth\else\Gin@nat@width\fi}
\def\maxheight{\ifdim\Gin@nat@height>\textheight0.8\textheight\else\Gin@nat@height\fi}
\makeatother
% Scale images if necessary, so that they will not overflow the page
% margins by default, and it is still possible to overwrite the defaults
% using explicit options in \includegraphics[width, height, ...]{}
\setkeys{Gin}{width=\maxwidth,height=\maxheight,keepaspectratio}

% Prevent slide breaks in the middle of a paragraph:
\widowpenalties 1 10000
\raggedbottom

\AtBeginPart{
  \let\insertpartnumber\relax
  \let\partname\relax
  \frame{\partpage}
}
\AtBeginSection{
  \ifbibliography
  \else
    \let\insertsectionnumber\relax
    \let\sectionname\relax
    \frame{\sectionpage}
  \fi
}
\AtBeginSubsection{
  \let\insertsubsectionnumber\relax
  \let\subsectionname\relax
  \frame{\subsectionpage}
}

\setlength{\parindent}{0pt}
\setlength{\parskip}{6pt plus 2pt minus 1pt}
\setlength{\emergencystretch}{3em}  % prevent overfull lines
\providecommand{\tightlist}{%
  \setlength{\itemsep}{0pt}\setlength{\parskip}{0pt}}
\setcounter{secnumdepth}{0}
\usepackage{graphicx}
\usepackage{amsmath,amsfonts,amssymb,amsthm,amscd, mathrsfs}
\setbeamertemplate{navigation symbols}{} %%removes bottom line
\setbeamertemplate{footline}[frame number]

\title{Class 9: Bayesian Hierarchical Models}
\author{Andrew Parnell\newline \texttt{andrew.parnell@mu.ie}
\newline \vspace{1cm} \newline \includegraphics[width=5cm]{MU_logo.jpg}}
\date{}

\begin{document}
\frame{\titlepage}

\begin{frame}[fragile]{Learning outcomes:}

\begin{itemize}
\tightlist
\item
  Start fitting hierarchical GLMs in \texttt{rstanarm}
\item
  Know some of the different versions of Hierarchical GLMs
\item
  Be able to expand and summarise fitted models
\end{itemize}

\end{frame}

\begin{frame}{From LMs to HGLMs}

\begin{itemize}
\tightlist
\item
  The Bayesian analogue of a mixed model is a \emph{hierarchical model}.
\item
  It's called a hierarchical model because the prior distributions come
  in layers, they depend on other parameters
\item
  Within this framework, we can borrow the ideas from the previous class
  to create hierarchical GLMs
\item
  We will go through four examples: binomial-logit, Poisson, and ordinal
  regression
\end{itemize}

\end{frame}

\begin{frame}[fragile]{Example 1: earnings data (again!)}

\begin{itemize}
\tightlist
\item
  Very easy to convert quickly from \texttt{lmer} to
  \texttt{stan\_lmer}:
\end{itemize}

\tiny

\begin{Shaded}
\begin{Highlighting}[]
\NormalTok{dat =}\StringTok{ }\KeywordTok{read.csv}\NormalTok{(}\StringTok{'data/earnings.csv'}\NormalTok{)}
\NormalTok{mod_}\DecValTok{1}\NormalTok{ =}\StringTok{ }\KeywordTok{stan_lmer}\NormalTok{(y }\OperatorTok{~}\StringTok{ }\NormalTok{x_centered }\OperatorTok{+}\StringTok{ }\NormalTok{(x_centered }\OperatorTok{|}\StringTok{ }\NormalTok{eth), }\DataTypeTok{data =}\NormalTok{ dat)}
\end{Highlighting}
\end{Shaded}

\begin{Shaded}
\begin{Highlighting}[]
\KeywordTok{round}\NormalTok{(}\KeywordTok{posterior_interval}\NormalTok{(mod_}\DecValTok{1}\NormalTok{), }\DecValTok{3}\NormalTok{)}
\end{Highlighting}
\end{Shaded}

\begin{verbatim}
##                                        5%   95%
## (Intercept)                         9.636 9.784
## x_centered                         -0.006 0.031
## b[(Intercept) eth:1]               -0.083 0.043
## b[x_centered eth:1]                -0.026 0.014
## b[(Intercept) eth:2]               -0.097 0.035
## b[x_centered eth:2]                -0.022 0.019
## b[(Intercept) eth:3]               -0.024 0.109
## b[x_centered eth:3]                -0.006 0.032
## b[(Intercept) eth:4]               -0.070 0.076
## b[x_centered eth:4]                -0.034 0.016
## sigma                               0.872 0.938
## Sigma[eth:(Intercept),(Intercept)]  0.000 0.016
## Sigma[eth:x_centered,(Intercept)]  -0.001 0.002
## Sigma[eth:x_centered,x_centered]    0.000 0.002
\end{verbatim}

\end{frame}

\begin{frame}[fragile]{Earnings model output}

\begin{Shaded}
\begin{Highlighting}[]
\KeywordTok{plot}\NormalTok{(mod_}\DecValTok{1}\NormalTok{)}
\end{Highlighting}
\end{Shaded}

\includegraphics{class_9_bhglms_files/figure-beamer/unnamed-chunk-3-1.pdf}

\end{frame}

\begin{frame}{Example 2: binomial-logit}

\begin{itemize}
\item
  Earlier we met the Binomial-logit model for binary data:
  \[y_i \sim Bin(1, p_i), logit(p_i) = \alpha + \beta (x_i - \bar{x})\]
  Here \(logit(p_i)\) is the link function equal to
  \(\log \left( \frac{p_i}{1-p_i} \right)\) and transforms the bounded
  probabilities into an unbounded space
\item
  If we have non-binary data we just change the likelihood:
  \[y_i \sim Bin(N_i, p_i), logit(p_i) = \alpha + \beta (x_i - \bar{x})\]
\item
  In a hierarchical version of this model, we vary the \emph{latent
  parameters} \(\alpha\) and \(\beta\) and give them prior distributions
\end{itemize}

\end{frame}

\begin{frame}[fragile]{The swiss willow tit data}

\tiny

\begin{Shaded}
\begin{Highlighting}[]
\NormalTok{swt =}\StringTok{ }\KeywordTok{read.csv}\NormalTok{(}\StringTok{'data/swt.csv'}\NormalTok{)}
\KeywordTok{head}\NormalTok{(swt)}
\end{Highlighting}
\end{Shaded}

\begin{verbatim}
##   rep.1 rep.2 rep.3 c.2 c.3 elev forest dur.1 day.2 day.3 length alt
## 1     0     0     0   0   0  420      3   240    58    73    6.2 Low
## 2     0     0     0   0   0  450     21   160    39    62    5.1 Low
## 3     0     0     0   0   0 1050     32   120    47    74    4.3 Med
## 4     0     0     0   0   0 1110     35   180    44    71    5.4 Med
## 5     0     0     0   0   0  510      2   210    56    73    3.6 Low
## 6     0     0     0   0   0  630     60   150    56    73    6.1 Low
\end{verbatim}

\end{frame}

\begin{frame}[fragile]{A hierarchical model}

\small
- Suppose we want to fit a model on the sum \(y_i =\)
\texttt{rep.1\ +\ rep.2\ +\ rep.3}:
\[y_i \sim Bin(N_i, p_i), logit(p_i) = \alpha_{\mbox{altitude}_i} + \beta_{\mbox{altitude}_i} (x_i- \bar{x})\]
where \(x_i\) is the percentage of forest cover

\begin{itemize}
\item
  What prior distributions should we use for \(\alpha\) and \(\beta\)?
\item
  Useful side note: A value of 10 on the logit scale leads to a
  probability of about 1, and a value of -10 leads to a probability of
  about 0 (you can test this by typing \texttt{inv.logit(10)}) so I
  wouldn't expect the value of \(logit(p_i)\) to ever get much bigger
  than 10 or smaller than -10
\item
  I have no idea whether we are more likely to find these birds in high
  percentage forest or low, so I'm happy to think that \(\beta\) might
  be around zero, and be positive or negative. Forest cover ranges from
  0 to 100 so that suggests that \(\beta\) is every likely to be bigger
  than 0.1 or smaller than -0.1. Perhaps \(\beta \sim N(0, 0.1^2)\) is a
  good prior
\item
  It looks to me like the intercept is very unlikely to be outside the
  range (-10, 10) so perhaps \(\alpha \sim N(0, 5^2)\) is appropriate
\end{itemize}

\end{frame}

\begin{frame}[fragile]{\texttt{rstanarm} code}

\begin{Shaded}
\begin{Highlighting}[]
\NormalTok{mod_}\DecValTok{2}\NormalTok{ =}\StringTok{ }\KeywordTok{stan_glmer}\NormalTok{(}\KeywordTok{cbind}\NormalTok{(y, N) }\OperatorTok{~}\StringTok{ }\NormalTok{forest }\OperatorTok{+}\StringTok{ }\NormalTok{(forest }\OperatorTok{|}\StringTok{ }\NormalTok{alt),}
                  \DataTypeTok{data =}\NormalTok{ swt,}
                  \DataTypeTok{family =} \KeywordTok{binomial}\NormalTok{(}\DataTypeTok{link =} \StringTok{'logit'}\NormalTok{),}
                  \DataTypeTok{prior =} \KeywordTok{normal}\NormalTok{(}\DecValTok{0}\NormalTok{, }\FloatTok{0.1}\NormalTok{),}
                  \DataTypeTok{prior_intercept =} \KeywordTok{normal}\NormalTok{(}\DecValTok{0}\NormalTok{, }\DecValTok{5}\NormalTok{))}
\end{Highlighting}
\end{Shaded}

\end{frame}

\begin{frame}[fragile]{Model summary 1}

\tiny

\begin{Shaded}
\begin{Highlighting}[]
\KeywordTok{posterior_interval}\NormalTok{(mod_}\DecValTok{2}\NormalTok{)}
\end{Highlighting}
\end{Shaded}

\begin{verbatim}
##                                               5%          95%
## (Intercept)                        -3.0547350656 -1.500090318
## forest                             -0.0053117955  0.006815317
## b[(Intercept) alt:High]            -0.1039793407  1.497335796
## b[forest alt:High]                  0.0131837840  0.033605020
## b[(Intercept) alt:Low]             -1.7271176612  0.020935355
## b[forest alt:Low]                  -0.0292270906  0.015086361
## b[(Intercept) alt:Med]             -0.3306274126  1.242020299
## b[forest alt:Med]                   0.0029792495  0.024435972
## Sigma[alt:(Intercept),(Intercept)]  0.0019747303  2.579853972
## Sigma[alt:forest,(Intercept)]      -0.0317051285  0.046441427
## Sigma[alt:forest,forest]            0.0001418906  0.025561489
\end{verbatim}

\end{frame}

\begin{frame}[fragile]{Model summary 2}

\begin{Shaded}
\begin{Highlighting}[]
\KeywordTok{plot}\NormalTok{(mod_}\DecValTok{2}\NormalTok{)}
\end{Highlighting}
\end{Shaded}

\includegraphics{class_9_bhglms_files/figure-beamer/unnamed-chunk-8-1.pdf}

\end{frame}

\begin{frame}[fragile]{Model fit - intercepts}

\begin{Shaded}
\begin{Highlighting}[]
\KeywordTok{mcmc_hist}\NormalTok{(}\KeywordTok{as.data.frame}\NormalTok{(mod_}\DecValTok{2}\NormalTok{), }
          \DataTypeTok{regex_pars =} \StringTok{'b}\CharTok{\textbackslash{}\textbackslash{}}\StringTok{[}\CharTok{\textbackslash{}\textbackslash{}}\StringTok{(Intercept'}\NormalTok{,}
          \DataTypeTok{facet_args =} \KeywordTok{list}\NormalTok{(}\DataTypeTok{nrow =} \DecValTok{3}\NormalTok{))}
\end{Highlighting}
\end{Shaded}

\includegraphics{class_9_bhglms_files/figure-beamer/unnamed-chunk-9-1.pdf}

\end{frame}

\begin{frame}[fragile]{Model fit - Slopes}

\begin{Shaded}
\begin{Highlighting}[]
\KeywordTok{mcmc_hist}\NormalTok{(}\KeywordTok{as.data.frame}\NormalTok{(mod_}\DecValTok{2}\NormalTok{), }\DataTypeTok{regex_pars =} \StringTok{'b}\CharTok{\textbackslash{}\textbackslash{}}\StringTok{[forest'}\NormalTok{)}
\end{Highlighting}
\end{Shaded}

\includegraphics{class_9_bhglms_files/figure-beamer/unnamed-chunk-10-1.pdf}

\end{frame}

\begin{frame}[fragile]{Model fit - posterior predictive check}

\begin{Shaded}
\begin{Highlighting}[]
\NormalTok{y_rep =}\StringTok{ }\KeywordTok{posterior_predict}\NormalTok{(mod_}\DecValTok{2}\NormalTok{)}
\NormalTok{y_rep_mean =}\StringTok{ }\KeywordTok{apply}\NormalTok{(y_rep, }\DecValTok{2}\NormalTok{, }\StringTok{'mean'}\NormalTok{)}
\KeywordTok{plot}\NormalTok{(swt}\OperatorTok{$}\NormalTok{y, y_rep_mean)}
\KeywordTok{abline}\NormalTok{(}\DataTypeTok{a =} \DecValTok{0}\NormalTok{, }\DataTypeTok{b =} \DecValTok{1}\NormalTok{, }\DataTypeTok{col =} \StringTok{'red'}\NormalTok{)}
\end{Highlighting}
\end{Shaded}

\includegraphics{class_9_bhglms_files/figure-beamer/unnamed-chunk-11-1.pdf}

\end{frame}

\begin{frame}[fragile]{Model fit - posterior predictive check 2}

\begin{Shaded}
\begin{Highlighting}[]
\KeywordTok{pp_check}\NormalTok{(mod_}\DecValTok{2}\NormalTok{)}
\end{Highlighting}
\end{Shaded}

\includegraphics{class_9_bhglms_files/figure-beamer/unnamed-chunk-12-1.pdf}

\end{frame}

\begin{frame}{Type 2: Poisson HGLMs}

\begin{itemize}
\item
  For a Poisson distribution there is no upper bound on the number of
  counts
\item
  We just change the likelihood (to Poisson) and the link function (to
  \(\log\)):
  \[y_i \sim Po(\lambda_i), \log(\lambda_i) = \alpha + \beta (x_i - \bar{x}))\]
\item
  We can now add our hierarchical layers into \(\alpha\) and \(\beta\),
  or\ldots{}
\item
  Another way we can add an extra layer is by giving \(\log(\lambda_i)\)
  a probability distribution rather than setting it to a value
\item
  This is a way of introducing \emph{over-dispersion}, i.e.~saying that
  the data are more variable than that expected by a standard Poisson
  distribution with our existing covariates
\end{itemize}

\end{frame}

\begin{frame}{An over-dispersed model}

\begin{itemize}
\item
  The over-dispersed model looks like:
  \[y_i \sim Po(\lambda_i), \log(\lambda_i) \sim N(\alpha + \beta (x_i - \bar{x}), \sigma^2)\]
  where \(\sigma\) is the over-dispersion parameter
\item
  We now need to estimate prior distributions for \(\alpha\), \(\beta\),
  and \(\sigma\)
\item
  We will use the horseshoe data again (see yesterday)
\end{itemize}

\end{frame}

\begin{frame}[fragile]{\texttt{rstanarm} code for OD Poisson}

\begin{Shaded}
\begin{Highlighting}[]
\NormalTok{horseshoe =}\StringTok{ }\KeywordTok{read.csv}\NormalTok{(}\StringTok{'data/horseshoe.csv'}\NormalTok{)}
\NormalTok{horseshoe}\OperatorTok{$}\NormalTok{obs <-}\StringTok{ }\DecValTok{1}\OperatorTok{:}\KeywordTok{nrow}\NormalTok{(horseshoe)}
\NormalTok{mod_}\DecValTok{3}\NormalTok{ =}\StringTok{ }\KeywordTok{stan_glmer}\NormalTok{(satell }\OperatorTok{~}\StringTok{ }\NormalTok{weight }\OperatorTok{+}\StringTok{ }\NormalTok{(}\DecValTok{1} \OperatorTok{|}\StringTok{ }\NormalTok{obs),}
              \DataTypeTok{family =}\NormalTok{ poisson, }\DataTypeTok{data =}\NormalTok{ horseshoe)}
\KeywordTok{mcmc_hist}\NormalTok{(}\KeywordTok{as.data.frame}\NormalTok{(mod_}\DecValTok{3}\NormalTok{), }\DataTypeTok{regex_pars =} \StringTok{'Sigma'}\NormalTok{)}
\end{Highlighting}
\end{Shaded}

\includegraphics{class_9_bhglms_files/figure-beamer/unnamed-chunk-13-1.pdf}

\end{frame}

\begin{frame}[fragile]{Posterior predictive check}

\begin{Shaded}
\begin{Highlighting}[]
\KeywordTok{pp_check}\NormalTok{(mod_}\DecValTok{3}\NormalTok{)}
\end{Highlighting}
\end{Shaded}

\includegraphics{class_9_bhglms_files/figure-beamer/unnamed-chunk-14-1.pdf}

\end{frame}

\begin{frame}{Notes about OD Poisson model}

\begin{itemize}
\item
  The way to think about OD models is via the data generating process.
  You could draw a DAG and think about how these processes might arise
\item
  We could compare this model to one without over dispersion via the PPC
  (or if time, WAIC, LOO, or cross validation).
\item
  In general, the parameter values (i.e.~the intercepts and slopes) tend
  to be more uncertain when you add in over dispersion
\end{itemize}

\end{frame}

\begin{frame}{Type 4: Ordinal data HGLMs}

\begin{itemize}
\tightlist
\item
  Often we have a response variable which is ordinal, e.g.~disagree,
  neutral, agree, etc
\item
  There are lots of different (and complicated) ways to model such data
\item
  Perhaps the easiest is to think of it as a hierarchical model with
  `cut-points' on a latent linear regression
\end{itemize}

\end{frame}

\begin{frame}{An ordinal model example}

\begin{itemize}
\tightlist
\item
  Suppose \(y_i = \{\mbox{disagree, neutral, agree} \}\) and we make it
  dependent on a latent continuous variable \(z_i\), so that :
\end{itemize}

\[y_i = \left\{ \begin{array}{ll} \mbox{agree} & \mbox{if } z_i> 0.5 \\
\mbox{neutral} & \mbox{if } -0.5 < z_i \le 0.5 \\
\mbox{disagree} & \mbox{if } z_i \le -0.5 \end{array} \right.\]

\begin{itemize}
\tightlist
\item
  We then give \(z_i\) a prior distribution, e.g.
  \(N(\beta_0 + \beta_1 x_i, \sigma^2)\)
\end{itemize}

\end{frame}

\begin{frame}[fragile]{Simulating some example data}

\begin{Shaded}
\begin{Highlighting}[]
\NormalTok{N =}\StringTok{ }\DecValTok{100}
\NormalTok{alpha =}\StringTok{ }\OperatorTok{-}\DecValTok{1}
\NormalTok{beta =}\StringTok{ }\FloatTok{0.2}
\NormalTok{sigma =}\StringTok{ }\FloatTok{0.51}
\KeywordTok{set.seed}\NormalTok{(}\DecValTok{123}\NormalTok{)}
\NormalTok{x =}\StringTok{ }\KeywordTok{runif}\NormalTok{(N, }\DecValTok{0}\NormalTok{, }\DecValTok{10}\NormalTok{)}
\NormalTok{cuts =}\StringTok{ }\KeywordTok{c}\NormalTok{(}\OperatorTok{-}\FloatTok{0.5}\NormalTok{, }\FloatTok{0.5}\NormalTok{)}
\NormalTok{z =}\StringTok{ }\KeywordTok{rnorm}\NormalTok{(N, alpha }\OperatorTok{+}\StringTok{ }\NormalTok{beta }\OperatorTok{*}\StringTok{ }\NormalTok{(x }\OperatorTok{-}\StringTok{ }\KeywordTok{mean}\NormalTok{(x)), sigma)}
\NormalTok{y =}\StringTok{ }\KeywordTok{findInterval}\NormalTok{(z, cuts)}
\NormalTok{dat =}\StringTok{ }\KeywordTok{data.frame}\NormalTok{(}\DataTypeTok{y =} \KeywordTok{as.factor}\NormalTok{(y),}
                 \DataTypeTok{x =}\NormalTok{ x)}
\end{Highlighting}
\end{Shaded}

\end{frame}

\begin{frame}[fragile]{Simulated data - plot}

\begin{Shaded}
\begin{Highlighting}[]
\KeywordTok{plot}\NormalTok{(x, z, }\DataTypeTok{col =}\NormalTok{ y }\OperatorTok{+}\StringTok{ }\DecValTok{1}\NormalTok{)}
\end{Highlighting}
\end{Shaded}

\includegraphics{class_9_bhglms_files/figure-beamer/unnamed-chunk-16-1.pdf}

\end{frame}

\begin{frame}[fragile]{Fitting in \texttt{rstanarm}}

\begin{Shaded}
\begin{Highlighting}[]
\NormalTok{mod_}\DecValTok{4}\NormalTok{ =}\StringTok{ }\KeywordTok{stan_polr}\NormalTok{(y }\OperatorTok{~}\StringTok{ }\NormalTok{x, }
                   \DataTypeTok{data =}\NormalTok{ dat,}
                   \DataTypeTok{prior =} \KeywordTok{R2}\NormalTok{(}\FloatTok{0.5}\NormalTok{, }\StringTok{"mean"}\NormalTok{),}
                   \DataTypeTok{prior_counts =} \KeywordTok{dirichlet}\NormalTok{(}\DecValTok{1}\NormalTok{))}
\end{Highlighting}
\end{Shaded}

\end{frame}

\begin{frame}[fragile]{Output}

\small

\begin{Shaded}
\begin{Highlighting}[]
\KeywordTok{mcmc_hist}\NormalTok{(}\KeywordTok{as.data.frame}\NormalTok{(mod_}\DecValTok{4}\NormalTok{))}
\end{Highlighting}
\end{Shaded}

\begin{verbatim}
## `stat_bin()` using `bins = 30`. Pick better value with `binwidth`.
\end{verbatim}

\includegraphics{class_9_bhglms_files/figure-beamer/unnamed-chunk-18-1.pdf}

\end{frame}

\begin{frame}{Summary}

\begin{itemize}
\tightlist
\item
  We have now seen a number of different types of hierarchical GLM
\item
  Many of the ideas of hierarchical linear models transfer over, but we
  can explore richer behaviour with hierarchical GLMs
\item
  These have all used the normal, binomial or Poisson distribution at
  the top level, and have allowed for over-dispersion, robustness, and
  ordinal data, to name just three
\end{itemize}

\end{frame}

\end{document}
